% Vietnamese translation for OI Public Library
% Source-File: IOI中国国家候选队论文/2003/方奇/方奇.doc
\documentclass[11pt]{article}
\usepackage{oipl-vn}

\title{Ứng dụng của phương pháp tô màu và phương pháp cấu tạo trên bàn cờ}
\author{Fang Qi}
\date{2003}

\begin{document}
\maketitle

\origtitle{Ứng dụng của phương pháp tô màu và phương pháp cấu tạo trên bàn cờ}
\sourcepath{IOI China candidate team papers / 2003 / Fang Qi / Fang Qi.doc}

\section{Khái niệm cơ bản}

\paragraph{Bàn cờ.}
Bàn cờ \(m\times n\) là hình chữ nhật gồm \(m\) hàng và \(n\) cột ô vuông.
Mỗi ô vuông được gọi là một ô của bàn cờ; ô ở hàng thứ \(i\), cột thứ \(j\)
được ký hiệu là \(a(i,j)\). Khi \(i+j\) là số lẻ, tương ứng số chẵn, ta gọi
\(a(i,j)\) là ô lẻ, tương ứng ô chẵn.

\paragraph{Phương pháp tô màu.}
Ta tô các ô của bàn cờ bằng nhiều màu khác nhau để đạt tác dụng phân loại.
Đặc biệt, kiểu tô hai màu đen trắng giống bàn cờ vua quốc tế được gọi là
\term{tô màu tự nhiên}.

\paragraph{Phương pháp cấu tạo.}
Phương pháp trực tiếp liệt kê một đối tượng toán học thỏa điều kiện, hoặc một
phản ví dụ làm cho kết luận được khẳng định hay phủ định; hoặc gián tiếp xây
dựng một quan hệ tương ứng để chuyển hóa bài toán theo nhu cầu, được gọi là
\term{phương pháp cấu tạo}.

\section{Phủ bàn cờ}

Phủ bàn cờ là dùng một số hình để phủ bàn cờ \(m\times n\). Mỗi hình dùng để
phủ cũng gồm một số ô, gọi là \term{hình phủ}. Ta quy ước hai hình phủ bất kỳ
không chồng lên nhau, và mỗi ô trong một hình phủ luôn trùng với một ô nào đó
trên bàn cờ.

Theo hiệu quả phủ, có thể chia thành phủ hoàn toàn, phủ bão hòa, phủ kín và
phủ khác nhau. Bài viết này chỉ xét vài trường hợp đơn giản. Trong đó, phủ
hoàn toàn nghĩa là tổng số ô của các hình phủ bằng tổng số ô của bàn cờ.

Theo loại hình phủ, có thể chia thành:
\begin{itemize}
  \item phủ đồng dạng: chỉ có một loại hình phủ;
  \item phủ dị dạng: có nhiều loại hình phủ.
\end{itemize}

\subsection{Phủ đồng dạng}

Cho \(m,n,k\). Hãy thử dùng một số hình chữ nhật \(1\times k\) để phủ bàn cờ
\(m\times n\).

\paragraph{Định lý 1.}
Bàn cờ \(m\times n\) tồn tại một cách phủ hoàn toàn bằng các hình chữ nhật
\(1\times k\) khi và chỉ khi
\[
  k\mid m \quad\text{hoặc}\quad k\mid n.
\]

\paragraph{Chứng minh.}
Chiều đủ là hiển nhiên và có thể chứng minh bằng cấu tạo. Nếu \(k\mid n\),
mỗi hàng được phủ đúng bằng \(n/k\) hình chữ nhật \(1\times k\). Trường hợp
\(k\mid m\) tương tự.

Với chiều cần, giả sử ngược lại rằng \(k\nmid m\) và \(k\nmid n\). Viết
\[
  m=m_1k+r,\quad 0<r<k,
\]
\[
  n=n_1k+s,\quad 0<s<k,
\]
và không mất tính tổng quát giả sử \(r\ge s\). Bản gốc minh họa bằng một sơ
đồ tô màu tuần hoàn theo chu kỳ \(k\): mỗi hình chữ nhật \(1\times k\), dù đặt
ngang hay dọc, đều phải phủ cùng một số ô của mỗi lớp màu. Nhưng vì hai phần
dư \(r,s\) đều khác \(0\), phần góc dư tạo ra một lớp màu có số lượng khác
với các lớp còn lại. Điều này mâu thuẫn với điều kiện cần của mọi cách phủ
hoàn toàn. Do đó ít nhất một trong hai số \(m,n\) phải chia hết cho \(k\).

Từ định lý trên, có thể giải triệt để bài toán phủ hoàn toàn bàn cờ
\(m\times n\) bằng hình chữ nhật \(p\times q\).

\paragraph{Định lý 2.}
Bàn cờ \(m\times n\) tồn tại một cách phủ hoàn toàn bằng các hình chữ nhật
\(p\times q\) khi và chỉ khi tồn tại cách chọn \(\{x,y\}=\{m,n\}\) sao cho
một trong các điều kiện sau đúng:
\begin{enumerate}
  \item \(p\mid x\) và \(q\mid y\);
  \item \(p\mid x\), \(q\mid x\), và tồn tại các số tự nhiên \(a,b\) sao cho
        \[
          y=ap+bq.
        \]
\end{enumerate}

\subsection{Phủ dị dạng}

\paragraph{Ví dụ 2.}
Xét bàn cờ \(m\times n\). Khi \(mn\) là số lẻ, hãy thử bỏ đi một ô, rồi dùng
một số hình chữ nhật \(1\times2\) phủ phần còn lại. Khi \(mn\) là số chẵn, hãy
thử bỏ đi hai ô, rồi dùng một số hình chữ nhật \(1\times2\) phủ phần còn lại.

\paragraph{Phân tích.}
Trước hết xét trường hợp \(mn\) lẻ. Tô màu tự nhiên bàn cờ. Dù đặt thế nào,
một hình chữ nhật \(1\times2\) luôn phủ đúng một ô đen và một ô trắng. Trên
bàn cờ ban đầu, số ô đen nhiều hơn số ô trắng \(1\), nên ô bị bỏ chỉ có thể
là một ô đen, tức một ô chẵn.

Ngược lại, giả sử ô bị bỏ là ô chẵn \(a(i,j)\). Bằng phương pháp cấu tạo, ta
có thể dựng được cách phủ khả thi. Bản gốc chia thành các trường hợp theo
tính chẵn lẻ của \(i\) và \(j\), rồi tách bàn cờ thành các dải rộng \(1\) để
đặt domino lần lượt trên từng dải. Khi \(i,j\) cùng lẻ hoặc cùng chẵn, các
dải được nối theo hai mẫu đối xứng; hai trường hợp còn lại suy ra tương tự.

Tiếp theo xét trường hợp \(mn\) chẵn. Từ phép tô màu tự nhiên, hai ô bị bỏ
phải khác màu, tức một ô lẻ và một ô chẵn; nếu không, số ô còn lại của hai
màu sẽ không bằng nhau.

Chiều ngược lại dùng cấu tạo. Ta dùng các đường đậm chia bàn cờ thành một
đường đi dạng dải rộng \(1\), sao cho từ bất kỳ ô nào cũng có thể đi qua mọi
ô đúng một lần rồi quay về điểm xuất phát. Khi bỏ đi một ô đen và một ô
trắng, chu trình này bị tách thành hai đường đi có số ô chẵn; mỗi đường đi
được phủ bằng cách ghép các ô liên tiếp thành các hình \(1\times2\).

\paragraph{Một nhận xét về tô màu.}
Trong các ví dụ trên, ta đều dùng điều kiện ``số ô của các màu phải bằng
nhau'' để suy ra mâu thuẫn. Tuy nhiên, có lúc chỉ xét điều kiện này vẫn chưa
đủ mạnh.

Chẳng hạn, với bàn cờ \(8\times8\), cần cắt đi ô nào để phần còn lại có thể
được phủ bằng \(21\) hình chữ nhật \(1\times3\)?

Ta dùng cách tô ba màu theo chu kỳ sao cho mỗi hình \(1\times3\) luôn phủ
đúng một ô của mỗi màu. Khi ấy số ô của ba màu lần lượt là
\[
  21,\quad 22,\quad 21.
\]
Vì vậy ô bị cắt phải thuộc màu xuất hiện \(22\) lần. Xét thêm đối xứng của
bàn cờ, chỉ các ô
\[
  a(3,3),\quad a(3,6),\quad a(6,3),\quad a(6,6)
\]
có thể thỏa yêu cầu.

\subsection{Tiểu kết}

Bài toán phủ là một chủ đề khá khó. Ở đây chỉ bàn đến vài trường hợp đơn giản
để minh họa ứng dụng của phương pháp tô màu và phương pháp cấu tạo. Cần bổ
sung rằng phương pháp tô màu có rất nhiều dạng. Xét đến khả năng mở rộng và
tính dễ thao tác, bài viết này chủ yếu nghiên cứu kiểu ``tô màu cách quãng'',
tức sự mở rộng của tô màu tự nhiên.

\section{Hành trình của quân mã}

Quy tắc đi của quân mã: từ một góc của hình chữ nhật \(2\times3\), quân mã
nhảy theo đường chéo đến góc đối diện. Các bài toán hành trình của quân mã
trên bàn cờ thường chia thành hai loại:
\begin{itemize}
  \item chuỗi Hamilton của quân mã;
  \item chu trình Hamilton của quân mã.
\end{itemize}

\subsection{Chuỗi Hamilton của quân mã}

Thông thường có một số phương pháp như sau:
\begin{itemize}
  \item tham lam: mỗi bước nhảy đến đỉnh có bậc nhỏ nhất;
  \item chia để trị: chia bàn cờ thành vài bàn cờ nhỏ, tìm chuỗi Hamilton
        trong từng phần rồi nối lại;
  \item viền mở rộng: trước hết tìm chuỗi Hamilton trong một bàn cờ nhỏ, rồi
        thêm viền bốn phía để sinh chuỗi Hamilton của bàn cờ lớn hơn.
\end{itemize}

Theo các phương pháp trên, không khó nhận được kết luận: bàn cờ \(n\times n\)
tồn tại chuỗi Hamilton của quân mã khi và chỉ khi
\[
  n>3.
\]

\subsection{Chu trình Hamilton của quân mã}

Xét bài toán tìm chu trình Hamilton của quân mã trên bàn cờ \(n\times n\).

Trước hết dùng tô màu tự nhiên để xét các trường hợp vô nghiệm. Dù quân mã đi
thế nào, nó cũng phải luân phiên giữa ô đen và ô trắng. Vì chu trình phải quay
về điểm xuất phát, và điểm xuất phát cùng màu với điểm kết thúc, số ô đi qua
của hai màu phải bằng nhau. Do đó khi \(n\) lẻ thì vô nghiệm. Do giới hạn về
kích thước, khi \(n<6\) cũng vô nghiệm.

Khi \(n\ge 6\) và \(n\) chẵn, có thể dùng phương pháp viền mở rộng để cấu
tạo. Giả sử ta đã tìm được một chu trình Hamilton trên bàn cờ
\((n-4)\times(n-4)\). Ta mở rộng ra bàn cờ \(n\times n\) bằng cách mở các
điểm nối ở bốn góc của hình chữ nhật bên trong, rồi nối chúng với bốn vòng
ngoài.

\paragraph{Trường hợp \(n\equiv2\pmod 4\).}
Mở chu trình bên trong tại bốn góc \(A,E,I,M\). Sau đó nối:
\[
\begin{array}{ll}
  1. & \text{vòng ngoài chứa } C,D \text{ với vòng trong tại } A,B,
       \text{ thành } A-C-\cdots-D-B,\\
  2. & \text{vòng ngoài chứa } G,H \text{ với vòng trong tại } E,F,
       \text{ thành } E-G-\cdots-H-F,\\
  3. & \text{vòng ngoài chứa } K,L \text{ với vòng trong tại } I,J,
       \text{ thành } I-K-\cdots-L-J,\\
  4. & \text{vòng ngoài chứa } O,P \text{ với vòng trong tại } M,N,
       \text{ thành } M-O-\cdots-P-N.
\end{array}
\]

Ở đây cần chú ý rằng chu trình của hình chữ nhật bên trong, dùng làm cơ sở,
trước hết phải thỏa: bước tiếp theo của \(A\) là \(B\), bước tiếp theo của
\(E\) là \(F\), bước tiếp theo của \(I\) là \(J\), và bước tiếp theo của \(M\)
là \(N\). Chỉ như vậy, hình chữ nhật mới cấu tạo được mới tiếp tục có thể làm
hình chữ nhật bên trong để mở rộng theo quy tắc trên. Bản gốc đưa ra một
nghiệm cơ sở thỏa yêu cầu cho \(N=6\).

\paragraph{Trường hợp \(n\equiv0\pmod 4\).}
Cách làm tương tự: mở chu trình bên trong tại bốn góc \(A,E,I,M\), rồi nối
bốn vòng ngoài theo đúng bốn phép nối
\[
  A-C-\cdots-D-B,\quad
  E-G-\cdots-H-F,\quad
  I-K-\cdots-L-J,\quad
  M-O-\cdots-P-N.
\]
Bản gốc đưa ra một nghiệm cơ sở tương ứng cho trường hợp này.

\paragraph{Một phỏng đoán.}
Với bàn cờ \(m\times n\), \(m\le n\), điều kiện cần và đủ để không tồn tại
chu trình Hamilton của quân mã là \(m,n\) thỏa một trong các điều kiện sau:
\begin{enumerate}
  \item \(m,n\) đều là số lẻ;
  \item \(m=1,2\) hoặc \(4\);
  \item \(m=3\) và \(n=4,6,8\).
\end{enumerate}

\section{Các ứng dụng khác}

\subsection{Ví dụ 5: Worm World, UVA}

Một con sâu bò trên một lưới \(N\times N\). Mỗi ô lưới có một số. Con sâu
không được đi qua cùng một số hai lần. Ban đầu, nó tùy ý chọn một ô làm điểm
xuất phát. Mỗi bước nó chỉ được bò theo phương ngang hoặc phương dọc, và không
được ra ngoài lưới. Hỏi con sâu phải di chuyển thế nào để đi qua nhiều ô nhất?
Bản gốc kèm một hình mẫu của đề.

\paragraph{Phân tích.}
Dùng phương pháp tô màu để tham lam ra một cận trên. Tô màu tự nhiên bàn cờ.
Gọi \(T_{\text{free}},T_{\text{black}},T_{\text{white}}\) lần lượt ghi nhận
số lượng của ba loại giá trị:
\begin{itemize}
  \item một số chỉ xuất hiện trên các ô trắng: tăng \(T_{\text{white}}\) thêm
        \(1\);
  \item một số chỉ xuất hiện trên các ô đen: tăng \(T_{\text{black}}\) thêm
        \(1\);
  \item một số xuất hiện trên cả ô trắng lẫn ô đen: tăng
        \(T_{\text{free}}\) thêm \(1\).
\end{itemize}

Giả sử \(T_{\text{white}}\le T_{\text{black}}\); nếu không thì đổi vai trò hai
màu. Khi đó cận trên của độ dài đường đi là
\[
L_{\max}=
\begin{cases}
  2\bigl(T_{\text{white}}+T_{\text{free}}\bigr)+1,
    & T_{\text{white}}+T_{\text{free}}<T_{\text{black}},\\
  T_{\text{white}}+T_{\text{free}}+T_{\text{black}},
    & T_{\text{white}}+T_{\text{free}}\ge T_{\text{black}}.
\end{cases}
\]

\section{Kết luận}

Bài toán tồn tại thường thích hợp với phương pháp tô màu; bài toán khả thi
thường thích hợp với phương pháp cấu tạo. Trong những bài toán lấy bàn cờ làm
mô hình, nếu biết phối hợp hai phương pháp này, ta thường đạt hiệu quả gấp
nhiều lần.

Một điểm bổ sung cho ví dụ bàn cờ \(8\times8\): nếu bỏ ô \((i,j)\) mà tồn tại
một phương án phủ, thì các ô đối xứng với nó,
\[
  (i,9-j),\quad (9-i,j),\quad (9-i,9-j),
\]
cũng phải cho những phương án phủ tương ứng. Vì vậy bốn điểm này phải cùng
thuộc lớp màu có thể bị xóa.

\end{document}
